% ====================================================================
% §18 전체 입력 인자와 피팅 준비 — 진행 요약
% ====================================================================
\section{전체 입력 인자와 피팅 준비 --- 진행 요약}\label{sec:inputs}
사용자 목표는 이 모든 인자를 입력으로 노출해 측정 데이터에 피팅하는 것이다. 곡선식의 모든 물리량이 사용자
조정 가능한 입력이며, 내부에 고정된 상수는 없다(전부 기본값$+$override --- 물리 상수 $R,F,k_B,h$ 와
수치 안전 가드 상수는 물리 인자가 아니므로 제외). 전 조정 인자와 기본값의 상세
목록은 부록~\ref{sec:appendix-code}(표~\ref{tab:inputs})에 정리한다 --- 어느 것을 고정하고 어느 것을
피팅할지는 사용자가 데이터로 정한다(본 장은 스코프를 정하지 않는다).

``고정$\cdot$피팅 스코프''는 자유롭게 잡을 수 있다 --- 예컨대
$\Delta H_\rxn$$\cdot$$\Delta S_\rxn$$\cdot$$Q$$\cdot$$w$ 를 풀고 $\Omega$$\cdot$$\Delta H_a$$\cdot$$|\dd V/\dd q|_{q_a}$
는 좁은 상하한으로 두는 식이다. 반응 열역학($\Delta H_\rxn\cdot\Delta S_\rxn$) 값과 DFT 활성화 값은 곧이
믿고 잠글 경계(상$\cdot$하한)가 아니라 \emph{출발점}이므로 소폭 자유도를 권한다.

\begin{keybox}
\textbf{본 장만으로 한 곡선 재현(6단계).}
(1) 표~\ref{tab:staging} 의 4 전이 파라미터 집합을 구성.
(2) 실험조건 $\sigma_d$, $|I|=$ c-rate$\,\cdot Q_\cell$, $T$, $R_n$ 설정 → 분극 $V_n$(\eqref{eq:n0map}\eqref{eq:vn}); 이후 평가는 $V_n$ 에서 점별.
(3) 전이마다 $U_j(T)$(\eqref{eq:Uj}) → 분기 중심 $U_j^{\,d}$(\eqref{eq:center}) → 폭 $w_j$(\eqref{eq:wbase}) → 평형 진행률 $\xi_\eq$(\eqref{eq:xieq}).
(4) 지연 길이: $\mathcal A$(\eqref{eq:Acut}) → $\chi_d,\Delta H_a^\eff$(\eqref{eq:chid}\eqref{eq:dHeff}) → $L_q$(\eqref{eq:Lqfull}) → $L_V$(\eqref{eq:LV}), 전이당 상수.
(5) peak 모양: 기억 적분(\eqref{eq:lag}$\cdot$\eqref{eq:reversal})의 $(\xi_\eq-\xi_\mathrm{lag})/L_V$(\eqref{eq:peakshape}); $L_V\to0$ 극한이 평형 종(\eqref{eq:tail-limit}\eqref{eq:eqpeak}).
(6) 용량 가중 합산(\eqref{eq:sum}); 출력이 입력 전위 $V_n$ 좌표에서 곧바로 나온다. 색인은 표~\ref{tab:nodemap}.
\end{keybox}

\subsection{전체 진행 요약}\label{sec:facade}
요약하면 --- 계산은 실험조건을 받아 방향 부호와 전류 크기로 먼저 환산하고(식~\eqref{eq:n0map}),
LCO 전극은 셀 라벨을 탈리튬화 부호로 자동 환산하며(식~\eqref{eq:lco-sigmaslot}), $|I|\to0$ 기준선은
식~\eqref{eq:eqpeak} 만 합산한 충방전 불변$\cdot$히스테리시스 미반영 곡선이다. 전체 진행은
표~\ref{tab:nodemap} 가 노드와 식으로 정리하고, 구현 대응표는
부록~\ref{sec:appendix-code} 에 둔다.

\begin{table}[t]
\centering
\renewcommand{\arraystretch}{1.3}
\caption{노드 $\leftrightarrow$ 식 매핑(전체 진행). 본 장 절 순서 $=$ 이 노드 순서.
구현 식별자 대응은 부록~\ref{sec:appendix-code}(표~\ref{tab:nodecode}).}
\label{tab:nodemap}
\begin{tabular}{l l l}
\toprule
노드 & 물리식 & 본 장 식 \\
\midrule
N0 & $\sigma_d$, $|I|=$ c-rate$\,\cdot Q_\cell$ & \eqref{eq:n0map} \\
N1 & $V_n=V_\app-\sigma_d|I|R_n$ & \eqref{eq:vn} \\
N2 & $U_j(T)=(-\Delta H+T\Delta S)/F$ & \eqref{eq:Uj} \\
N3 & $\Delta U_j^\hys$, $U_j^{\,d}$ & \eqref{eq:dUhys},\eqref{eq:Ubranch} \\
N4 & $w_j=n_jRT/F$ (이중지위) & \eqref{eq:wbase} \\
N5 & $\xi_\eq=\mathrm{logistic}[\sigma_d(V-U^d)/w]$ & \eqref{eq:xieq} \\
N6 & $Q_j\xi_\eq(1-\xi_\eq)/w$ & \eqref{eq:eqpeak} \\
\multicolumn{3}{l}{\emph{\quad(N6 확장: 두-상 broadening$\cdot$현상학적 $w_j$$\cdot$역산 금지 = \S\ref{sec:broadening}; 새 N-노드 아님, 모델 항 0)}}\\
N7 & $\mathcal A,\chi_d,\Delta H_a^\eff,L_q,L_V$ & \eqref{eq:Acut}--\eqref{eq:LV} \\
N8 & $(\xi_\eq-\xi_\mathrm{lag})/L_V$, 역전 & \eqref{eq:peakshape},\eqref{eq:reversal} \\
N9 & $C_\bg+\sum_jQ_j[\cdot]$ & \eqref{eq:sum} \\
\midrule
\multicolumn{3}{l}{\emph{— LCO 양극 추가(같은 골격, 파라미터+1항) —}}\\
N0$'$ & LCO 전이 초기값(3 전이) & 표~\ref{tab:lco-staging} \\
\multicolumn{3}{l}{\emph{\quad(N0$'$ 보강: 셀 라벨$\to$탈리튬화 부호 자동 환산 = 식~\eqref{eq:lco-sigmaslot}$\cdot$\S\ref{sec:facade})}}\\
N2$'$ & $\partial U_j/\partial T=\Delta S^\mathrm{cat}/F$ & \eqref{eq:lco-dUdT} \\
N3$'$ & $\Delta U_j^{\hys,\mathrm{cat}}$, $U_j^{\,d,\mathrm{cat}}$ & \eqref{eq:lco-dUhys},\eqref{eq:lco-Ubranch} \\
N5$+$ & $\Delta S_{e,j}^{\,\mathrm{mol}}=N_A\tfrac{\pi^2}{3}k_B^2T\,\partial g/\partial x<0$ (삽입, 몰당) & \eqref{eq:dSemolar},\eqref{eq:ggate} \\
N9$'$ & $\Delta S^\mathrm{cat}=$config$+$vib$+$elec & \eqref{eq:lco-decomp} \\
\bottomrule
\end{tabular}
\end{table}

여기까지가 본 장(흑연)의 결과 사슬이다 --- 표의 LCO 행(N$'$ 계열)은 Chapter 2 가 같은 forward 골격을
재사용하는 지점의 통합 대응 참조다. 부록 A 가 부호 사슬을
전수 검산하고, 부록 B 가 구현과의 1:1 대응을 역방향으로 조회한다.

% 자산: [B2-127] [B2-128] [B2-129]
